This study examines whether confident wording increases the scores assigned to incorrect mathematical reasoning steps. We use 50 selected MATH-500 questions to construct paired correct and incorrect targets, with each error formed by increasing an integer result by one. A frozen 1.5-billion-parameter Skywork process reward model evaluates each target under baseline, neutral, and confident wording, producing 300 scores. The neutral and confident variants have equal token counts. Relative to neutral wording, the estimated mean change in incorrect-step scores is \WrongMean{} (95\% paired-bootstrap interval: \WrongCI{}). Correct-step scores increase by \CorrectMean{} on average, while the estimated mean reduction in the correct--incorrect score gap is \ReductionMean{} (95\% paired-bootstrap interval: \ReductionCI{}). Both intervals include zero, leaving the direction of these two mean effects uncertain. Post hoc comparisons show narrower mean gaps under both added phrases than under the unprefixed baseline. Correct targets retain higher scores than their paired incorrect targets in all three wording conditions.
